\documentclass[11pt,a4paper]{article}
\usepackage[T1]{fontenc}
\usepackage[utf8]{inputenc}
\usepackage{lmodern}
\usepackage[margin=28mm]{geometry}
\usepackage{amsmath,amssymb,amsthm,mathtools}
\usepackage{booktabs,longtable,array}
\usepackage{enumitem,microtype}
\usepackage[hidelinks]{hyperref}
\usepackage{url}
\ifdefined\pdftrailerid\pdftrailerid{}\fi
\ifdefined\pdfsuppressptexinfo\pdfsuppressptexinfo=-1\fi
\providecommand{\SourceCommit}{SOURCE-NOT-BUILT}
\newtheorem{theorem}{Theorem}[section]
\newtheorem{lemma}[theorem]{Lemma}
\newtheorem{proposition}[theorem]{Proposition}
\newtheorem{corollary}[theorem]{Corollary}
\theoremstyle{definition}\newtheorem{definition}[theorem]{Definition}
\theoremstyle{remark}\newtheorem{remark}[theorem]{Remark}
\newcommand{\R}{\mathbb R}
\newcommand{\C}{\mathbb C}
\newcommand{\Z}{\mathbb Z}
\newcommand{\Exp}{\mathbb E}
\newcommand{\En}{\mathcal E}
\newcommand{\UU}{\mathcal U}
\newcommand{\Id}{I}
\newcommand{\ip}[2]{\langle #1,#2\rangle}
\newcommand{\norm}[1]{\lVert #1\rVert}
\DeclareMathOperator{\im}{im}
\DeclareMathOperator{\rank}{rank}
\DeclareMathOperator{\tr}{tr}
\DeclareMathOperator{\diag}{diag}
\title{Canonical rectangle observations for matrix gradients:\\
fixed weak spaces and checkerboard counterexamples}
\author{Lluis Eriksson\\\small Independent researcher}
\date{Source preparation: 11 November 2026\\\small Revision 250}
\hypersetup{pdftitle={Canonical rectangle observations for matrix gradients: fixed weak spaces and checkerboard counterexamples},pdfauthor={Lluis Eriksson}}
\begin{document}
\maketitle
\begin{abstract}
We study finite positive matrix-weighted gradients on a periodic lattice, with
observations restricted to normalized rectangular means in a fixed internal
basis. For a spatially constant rank-$q$ weak projector in $\R^g$ satisfying
$P_{\alpha\alpha}\leq\Lambda q/g$, the minimum actual restoring rank on the
flat coexact domain has order $q n^D+g n^{D-2}$, uniformly in $q$.
An explicit choice of original-channel rectangle partitions attains this order.
For this constructed map, the least quadratic coefficient has order
$c\min(g/q,n^2)$. The construction uses a stated theorem of
Marcus--Spielman--Srivastava, augmented to control both Gram matrices and
cardinalities, and derived finite Neumann estimates.
Spatial constancy is substantive. On an exact positive two-channel checkerboard
with the same local spectra and fixed basis, full output counts are
$O(n\log n)$ in dimension three and $O(n^2)$ in dimension four, with finite
size-dependent penalties. We derive the point estimates, actual strong graphs,
integer meshes and original-rectangle realization. A separate plane construction
works at every fixed nonzero angle between the two weak lines.
These are finite observation-class results; optimal checkerboard rank, practical
conditioning and historical priority remain unresolved.
\end{abstract}
\tableofcontents
\section{The observation class and the main results}\label{sec:setup}
The problem is to restore a fixed coercivity target using few permitted
observations. Two distinctions matter: the dimension of a displayed output
space versus its actual rank on the stated domain, and a finite penalty versus
a penalty bounded independently of the lattice size. We keep both explicit.
The proofs concern one positive family of finite operators, rather than a
uniform theorem for arbitrary matrix profiles.

\subsection{Actual finite gradients and the coexact domain}
Fix integers $D\geq3$, $n\geq2$ and $g\geq1$, and write
\[
 X=(\Z/n\Z)^D,\qquad V=n^D,\qquad
 \mathcal H_g=(\R^D\otimes\R^g)^X.
\]
All unqualified inner products and norms use counting measure. Put
$\partial_i f(x)=f(x+e_i)-f(x)$, retaining every positive forward label,
including period-two multiplicities. For $z\in\mathcal H_g$ let
\[
 \En(z)=\sum_{x\in X}\sum_{i,k=1}^D\sum_{\alpha=1}^g
       |\partial_i z_{k,\alpha}(x)|^2.
\]
For a real one-form $v$ define
$(d_1v)_{ij}=\partial_i v_j-\partial_j v_i$ for $i<j$.
The adjoint uses the counting inner products, with each unordered pair $i<j$
counted once. Our domain is
\begin{equation}\label{eq:domain}
 \UU_0=\im d_1^*,\qquad \UU=\UU_0^g\subset\mathcal H_g,
 \qquad d_0:=\dim_\R\UU_0=(D-1)(V-1).
\end{equation}
The dimension identity is proved in Lemma~\ref{lem:fourier}. No artificial
componentwise replacement of this domain is used.

For a positive symmetric $g\times g$ field $B(x)$ set
\begin{equation}\label{eq:energy}
 Q_B(z)=\sum_{x\in X}\sum_{i,k=1}^D
       (\partial_i z_k(x))^T B(x)\partial_i z_k(x).
\end{equation}
The coefficient is evaluated at the actual forward \emph{source} $x$.
The full gradient-stack profile is $\mathcal A(x)=I_{D^2}\otimes B(x)$;
there are $D^2g$ local eigenvalues, while the field has $h=Dg$ components.
This notation does not assert an identity with an interacting or continuum
operator.

\subsection{Canonical observations and normalization}
For each original internal channel $\alpha$, choose a partition
$\mathcal P_\alpha$ of $\{0,\ldots,n-1\}^D$ into nonempty integer rectangles
$S=\prod_i\{a_i,\ldots,b_i\}$. Different channels may have different partitions;
within a channel the partition is shared across all spatial components $k$.
The canonical observation map has rows
\begin{equation}\label{eq:canonical}
 (Cz)_{\alpha,k,S}=V^{-1/2}|S|^{-1/2}
             \sum_{x\in S}z_{k,\alpha}(x),\qquad
 N_{\rm out}=D\sum_\alpha|\mathcal P_\alpha|.
\end{equation}
These are unrotated, fixed-weight component means. No mixed internal row or
subsequent dimension-reducing channel mixture is part of this class.
Let $\mathcal R$ be the counting orthogonal block projection replacing each
component by its arithmetic mean on its channel's rectangles. Then
\begin{equation}\label{eq:RC}
 \norm{Cz}^2=V^{-1}\norm{\mathcal Rz}^2.
\end{equation}
The \emph{actual rank} is $r(C)=\rank(C|_\UU)$, which need not be $N_{\rm out}$.
We call $C$ restoring at target $c>0$ if some finite $K\geq0$ gives
\begin{equation}\label{eq:restoring}
 V^{-1}Q_B(z)+K\norm{Cz}^2\geq cV^{-1}\norm z^2
 \quad(z\in\UU).
\end{equation}
We separately state when the same inequality holds on all of $\mathcal H_g$.
On these \emph{same} rows the counting-energy convention uses coefficient $VK$.
Arbitrary row rescaling would change the penalty question and is excluded.
The projection $\mathcal R$ is not assumed to preserve $\UU$.

\subsection{Sharp order for a fixed weak space}
Let $1\leq q\leq g$, $\Lambda\geq1$, and let $P$ be a spatially constant
rank-$q$ orthogonal projector satisfying
\begin{equation}\label{eq:coherence}
 P_{\alpha\alpha}\leq\Lambda q/g\quad(1\leq\alpha\leq g).
\end{equation}
With $c>0$ fixed put
\begin{equation}\label{eq:profile}
 a=\frac{c}{16D},\qquad H=cn^2,\qquad
 B=HI_g-(H-a)P,\qquad b=\min\{g/q,n^2\}.
\end{equation}
There is no lower bound on nonzero diagonal entries of $P$ and no prescribed
well-conditioned subset of channels.

\begin{theorem}[Fixed weak space]\label{thm:fixed}
For \eqref{eq:profile}--\eqref{eq:coherence} there are coefficient-only canonical
partitions such that
\begin{equation}\label{eq:fixedupper}
 N_{\rm out}\leq L_{D,\Lambda}(qV+gV/n^2),\qquad
 Q_B(z)/V+32cb\norm{Cz}^2\geq c\norm z^2/V
 \quad(z\in\mathcal H_g).
\end{equation}
One valid constant is
\[
 C_0=2^{24},\quad
 A_D=\frac{D2^D(17/16)}{2^{D-2}-1},\quad
 L_{D,\Lambda}=A_D+4352D+\frac{17}{4}DC_0\Lambda.
\]
Every restoring internally channel-separable linear map, even beyond canonical
rectangles, has actual rank
\begin{equation}\label{eq:ranklower}
 r\geq\max\left\{qd_0,
 \frac{gd_0}{2\Lambda(1+8Dn^2)}\right\}.
\end{equation}
This rank bound only uses the restoring inequality at zero observation; it
also holds for any finite penalty $\Phi(Cz)$ with $\Phi(0)=0$ in place of the
quadratic penalty. Consequently the minimum canonical restoring actual rank
has order $qV+gV/n^2$, with constants depending only on $D,\Lambda$.
For the particular constructed canonical map, its least quadratic coefficient
on either $\UU$ or all fields satisfies
\begin{equation}\label{eq:penaltyorder}
 \frac{3(D-1)}{16\Lambda L_{D,\Lambda}}cb
 \leq K_{\min}\leq32cb.
\end{equation}
All constants are independent of $q,g,n,c$.
\end{theorem}
Here channel-separable means a direct sum of linear maps, one acting on each
original channel's entire $D$-component field. No lower bound for internally
mixed maps is asserted. The coefficient order in \eqref{eq:penaltyorder}
concerns a fixed constructed map, not an upper bound for every map of that rank.
The construction is a finite existence prescription; neither efficient search
nor a practical condition number is established.

\subsection{Spatial variation with the same local spectra}
For even $n$ let $p(x)=\sum_i x_i\bmod2$, which is well-defined on the torus.
For $g=2,q=1,\Lambda=2$, set
\begin{equation}\label{eq:checker}
 P(x)=\begin{cases}e_1e_1^T,&p(x)=0,\\e_2e_2^T,&p(x)=1,\end{cases}
 \qquad B(x)=HI_2-(H-a)P(x),
\end{equation}
with the same $a,H$ from \eqref{eq:profile}. Only spatial constancy of $P$
has been removed; local eigenvalues, positivity, upper coherence, observed
basis, canonical weights and actual domain are unchanged.

\begin{theorem}[Orthogonal checkerboard]\label{thm:checker}
For \eqref{eq:checker} and every even $n\geq2$:
\begin{enumerate}[label=(\roman*),leftmargin=*]
\item For every $D\geq3$ a canonical map with
$N_{\rm out}=2D(n^{D-1}+1)$ restores $c$ on all fields with a finite coefficient.
\item For $D=3$ a canonical map with
$N_{\rm out}\leq2^{19}n(1+\ln n)$ restores $c$ on all fields.
\item For $D=4$ a canonical map with
$N_{\rm out}\leq2^{25}n^2$ restores $c$ on all fields.
\end{enumerate}
In (ii)--(iii), the non-singleton construction has the sufficient normalized
coefficient
\begin{equation}\label{eq:pointK}
 K=M+\frac{2M^2}{3c},\qquad M=4DH+c;
\end{equation}
the explicit singleton fallbacks use $K=c$. Actual rank on $\UU_0^2$ is at
most the corresponding full count. These are upper bounds, not optimality,
matching lower bounds, logarithmic necessity or uniform-penalty statements.
\end{theorem}
In particular a universal per-profile $\Omega(V)$ rank lower cannot be extended
from constant $P$ merely by retaining the same positive local spectra and
pointwise coherence. This does not refute a worst-case lower over a larger
class containing constant profiles, or a lower with an additional uniform
penalty restriction.

\begin{theorem}[Fixed-angle plane construction]\label{thm:angle}
Replace the two projectors in \eqref{eq:checker} by any two fixed distinct weak
lines with principal angle $\theta\in(0,\pi/2]$, in the original observed basis.
Let $\kappa=1-|\cos\theta|$ and $m_\theta=\lceil4/\sqrt\kappa\rceil$.
For every $D\geq3$ and even $n\geq2$ there is an all-field restoring canonical
map with
\[
 N_{\rm out}\leq2Dm_\theta(n^{D-1}+1)
\]
and a finite penalty. When $n\geq m_\theta$, \eqref{eq:pointK} suffices;
otherwise the singleton fallback has $K=c$.
\end{theorem}
For fixed $\theta,D$ this is $o(V)$, with constants allowed to diverge as
$\theta\downarrow0$. We prove no $n$-dependent-angle order and no extension of
the point-anchor orders in Theorem~\ref{thm:checker}(ii)--(iii) to nonorthogonal
lines. That separate question is outside the present theorem list.

\subsection{The unchanged capped budget and proof roadmap}
To fix the normalization inherited by these examples, define for positive
local gradient-stack eigenvalues $\lambda_j(x)$
\begin{equation}\label{eq:cap}
 \alpha_D=\left(\frac34\right)^D\frac D2
       \left(\frac D{D-2}\right)^{D-1},\qquad
 F_c(\mathcal A)=\sum_{x,j}\min\{1,\alpha_D(c/\lambda_j(x))^{D/2}\}.
\end{equation}
For the fixed-space family, $a$ has multiplicity $qD^2$ and $H$ has
multiplicity $(g-q)D^2$ per site. Its weak cap equals one, so
\begin{equation}\label{eq:fixedcap}
 F_c(\mathcal A)=D^2V[q+(g-q)\min\{1,\alpha_Dn^{-D}\}]
 \leq D^2(qV+\alpha_Dg).
\end{equation}
Indeed $\alpha_D(16D)^{D/2}\geq(3\sqrt D)^D>1$.
Thus $h+F_c(\mathcal A)=O_D(qV+g)$; we do not identify it with the optimal
rank expression $qV+gV/n^2$. For both checkerboards the same formula with
$q=1,g=2$ is unchanged. In particular
\[
 \alpha_3=729/128,\quad \alpha_4=81/16,\qquad
 F_c=D^2V[1+\min\{1,\alpha_Dn^{-D}\}],\quad h=2D.
\]
Every cap and multiplicity is retained.

Section~\ref{sec:finite} derives all elementary finite spectral and projection
facts. Section~\ref{sec:fixed} proves Theorem~\ref{thm:fixed}; its sole
non-elementary input is the explicitly stated primary MSS theorem.
Sections~\ref{sec:checkerproof}--\ref{sec:points} prove the actual strong-graph,
plane and point constructions. Scope, limited primary comparisons and exact
source/review context are separated in Section~\ref{sec:context}.

\section{Finite spectral and projection facts}\label{sec:finite}
This section derives the local estimates used later; a regional gap is not
an additional hypothesis.

\begin{lemma}[Actual Fourier domain and periodic gap]\label{lem:fourier}
The dimension in \eqref{eq:domain} is $d_0=(D-1)(V-1)$. Every element of
$\UU_0$ has mean zero. On every componentwise mean-zero field,
\[
 \En(f)\geq\frac{16}{n^2}\norm f^2,
 \qquad\text{while on every field }\En(f)\leq4D\norm f^2.
\]
\end{lemma}
\begin{proof}
Use the counting-orthonormal characters
$V^{-1/2}\exp(2\pi i\ell\cdot x/n)$ with
$\ell\in\{0,\ldots,n-1\}^D$. The forward difference multiplier is
$\xi_i(\ell)=\exp(2\pi i\ell_i/n)-1$. Its energy multiplier is
$4\sum_i\sin^2(\pi\ell_i/n)$, which is at most $4D$.
At a nonzero frequency at least one term is at least
$4\sin^2(\pi/n)\geq16/n^2$, by concavity of sine on $[0,\pi/2]$;
the case $n=2$ gives equality and retains the actual forward labels.

The curl symbol is $v\mapsto\xi\wedge v$. When $\xi\ne0$, its kernel is
$\C\xi$, so the adjoint image has complex dimension $D-1$; at zero frequency
that image is zero. Conjugate frequencies pair to give twice their complex
dimension as a real dimension. At self-conjugate frequencies each $\xi_i$ is
real ($0$ or $-2$), and the real adjoint image has dimension $D-1$ directly.
Summing gives $(D-1)(V-1)$, without an omitted period-two mode. The zero
frequency is absent, proving the mean-zero assertion. The upper energy bound
also follows from the counting isometry of every shift.
\end{proof}

\begin{lemma}[Finite counting Neumann basis]\label{lem:neumann}
On an $\ell$-vertex path, $\ell\geq2$, the counting-orthonormal vectors are
\[
 \phi_0(t)=\ell^{-1/2},\qquad
 \phi_j(t)=\sqrt{2/\ell}\cos\frac{\pi j(t+1/2)}\ell
 \quad(1\leq j<\ell),
\]
with eigenvalues $4\sin^2(\pi j/(2\ell))$. A singleton has only its constant
mode. On any product rectangle with side lengths at most $L$, the counting
variance is at most $L^2/4$ times its internal edge energy, also for vector
fields; singleton factors contribute zero variance.
\end{lemma}
\begin{proof}
For $\theta=\pi j/\ell$, interior second differences of
$\cos(\theta(t+1/2))$ multiply it by $2(1-\cos\theta)$. The virtual endpoint
values $f(-1)=f(0)$ and $f(\ell)=f(\ell-1)$ give precisely endpoint degree
one. For an integer $1\leq r<2\ell$,
\[
 \sum_{t=0}^{\ell-1}\cos\frac{\pi r(t+1/2)}\ell
 =\frac{\sin(\pi r)}{2\sin(\pi r/(2\ell))}=0.
\]
The denominator is nonzero in this range. Product-to-sum, with the $r=0$
sum equal to $\ell$, proves the stated orthogonality and normalization.
The needed sums $j+j'$ are at most $2\ell-2$. There are $\ell$ orthonormal
vectors, so this is a complete basis.

Tensor products diagonalize the product-path energy. Since
$\sin t\geq2t/\pi$ on $[0,\pi/2]$, every nonconstant product mode has
eigenvalue at least $4/L^2$. Removing the constant mode gives the variance
bound by expansion in this finite basis. If the entire rectangle is a
singleton there is no positive eigenvalue to invoke and the variance is zero.
Sum the scalar inequalities for vector fields. No dimension factor is lost.
\end{proof}

For an integer $j\geq0$ partition each coordinate chart into the nonempty
intervals
\[
 \{\lfloor rn/2^j\rfloor,\ldots,\lfloor(r+1)n/2^j\rfloor-1\},
 \qquad0\leq r<2^j,
\]
and take their products. Call the partition $\mathcal P_j$ and its counting
averaging projection $\Pi_j$. Every coarse endpoint is an even-index fine
endpoint, so the partitions are nested after empty intervals are omitted.
There are at most $2^{Dj}$ cells and their sides have length at most
$L_j=\lceil n/2^j\rceil$. In particular $\Pi_0$ is global averaging.

\begin{lemma}[Actual spatial pieces]\label{lem:pieces}
Fix $k\geq0$. The operators
\[
 \Delta_j=\Pi_{j+1}-\Pi_j\quad(0\leq j<k),\qquad
 \Delta_k=I-\Pi_k
\]
are mutually orthogonal projections with sum $I-\Pi_0$. For every field $f$,
\begin{equation}\label{eq:piecegap}
 \norm{\Delta_jf}^2\leq\frac{L_j^2}{4}\En(f)
 \quad(0\leq j\leq k).
\end{equation}
For $j<k$, $\Delta_j$ is a subprojection of $\Pi_{j+1}$; the tail is a
subprojection of the singleton observation $I$.
\end{lemma}
\begin{proof}
Nested averaging projections satisfy $\Pi_i\Pi_j=\Pi_{\min(i,j)}$.
Their differences and the displayed tail are therefore orthogonal projections
and telescope. Apply Lemma~\ref{lem:neumann} on every parent cell of
$\mathcal P_j$ and sum. Its internal edges are a subset of the original
labelled periodic edges, counted once. Other edges add nonnegative energy.
This gives $\norm{(I-\Pi_j)f}^2\leq L_j^2\En(f)/4$.
Each $\Delta_j$ is a subprojection of $I-\Pi_j$, proving
\eqref{eq:piecegap}, including the single tail. The subprojection statements
follow from the same identities. Wrap edges are never inserted into a
Neumann cell as fictitious interior edges.
\end{proof}

\begin{lemma}[Finite kernel-to-all-field completion]\label{lem:completion}
Let $L$ be a positive semidefinite operator on a finite real Hilbert space,
let $R$ be an orthogonal projection, and suppose
$\ip u{Lu}\geq(c+\delta)\norm u^2$ on $\ker R$, with $c,\delta>0$.
If $\norm{L-cI}\leq M$, then
\begin{equation}\label{eq:completion}
 \ip z{Lz}+\left(M+\frac{2M^2}{\delta}\right)\norm{Rz}^2
 \geq c\norm z^2\quad\text{for all }z.
\end{equation}
For \eqref{eq:energy} with $0<B(x)\leq HI$, one may take $M=4DH+c$.
\end{lemma}
\begin{proof}
Write $z=u+v$, $u=(I-R)z$, $v=Rz$, and $F=L-cI$.
Then $\ip u{Fu}\geq\delta\norm u^2$, $\ip v{Fv}\geq-M\norm v^2$, and
\[
 2|\ip u{Fv}|\leq2M\norm u\norm v
 \leq\frac\delta2\norm u^2+\frac{2M^2}\delta\norm v^2.
\]
After adding the displayed penalty a nonnegative $\delta\norm u^2/2$ remains.
The final operator bound follows from Lemma~\ref{lem:fourier}'s upper energy
bound and $B(x)\leq HI$. Nothing requires $R$ to preserve a later restricted
domain. Constants must still be controlled by $R$ or a kernel margin; they
are not silently discarded by this lemma.
\end{proof}

\begin{lemma}[Original-coordinate singleton isolation]\label{lem:isolation}
Given $A$ distinct points in the original $D$-dimensional chart, there is a
partition into at most $1+2DA$ nonempty integer rectangles in which every
specified point is a singleton.
\end{lemma}
\begin{proof}
Start with the whole chart. In the rectangle containing the next point,
cut before and after its first coordinate and retain the two possible side
rectangles. Continue only within the middle slice, through the remaining
coordinates. At most $2D$ side rectangles and the desired singleton replace
one rectangle, a net increase at most $2D$. Omit empty pieces. Disjoint
coverage follows from the first coordinate at which a point differs from
the anchor. Earlier singleton anchors lie in other rectangles and remain
unchanged; an already-singleton point costs zero. Iteration proves the count.
This uses no global tensor-grid propagation of cuts or point separation.
\end{proof}

\section{A fixed coherent weak space}\label{sec:fixed}
We prove Theorem~\ref{thm:fixed}. Write $\mathsf U\in\R^{g\times q}$ for an
orthonormal basis of $\im P$, so $\mathsf U\mathsf U^T=P$. One deterministic
choice is Gram--Schmidt on the first independent columns of $P$, in original
index order. If $u_\alpha\in\R^q$ is its $\alpha$th row, then
\begin{equation}\label{eq:frame}
 \sum_{\alpha=1}^g u_\alpha u_\alpha^T=I_q,\qquad
 \norm{u_\alpha}^2=P_{\alpha\alpha}\leq d:=\Lambda q/g.
\end{equation}
These coordinates serve the proof and channel selection. The observed fields
remain the original $z_{k,\alpha}$.

\subsection{The primary MSS input and an added count coordinate}
We use precisely the following classical theorem of
Marcus--Spielman--Srivastava \cite[Theorem 1.4, p.~3]{MSS}, in its version-pinned
arXiv statement.
\begin{theorem}[MSS finite-support theorem]\label{thm:MSS}
Let independent finitely supported random vectors $v_1,\ldots,v_m\in\C^N$
satisfy $\sum_i\Exp(v_iv_i^*)=I_N$ and
$\Exp\norm{v_i}^2\leq\varepsilon$ for every $i$, where $\varepsilon>0$.
With positive probability,
\[
 \norm{\sum_i v_iv_i^*}\leq(1+\sqrt\varepsilon)^2.
\]
\end{theorem}
We use this published theorem, not a proof or computational verification of it.
The simultaneous cardinality assertion needed next is not simply assumed
from a frame-partition corollary.

\begin{lemma}[Simultaneous Gram and count halving]\label{lem:halving}
Suppose $\sum_{i\in I}w_iw_i^T=I_q$, $m=|I|$, and
$\delta=\max_i\norm{w_i}^2+1/m$. There is a partition $I=S_1\sqcup S_2$
such that, for both $j=1,2$,
\begin{equation}\label{eq:half}
 (\tfrac12-e)I_q\leq\sum_{i\in S_j}w_iw_i^T\leq(\tfrac12+e)I_q,
 \qquad \tfrac12-e\leq |S_j|/m\leq\tfrac12+e,
 \quad e=\sqrt{2\delta}+\delta.
\end{equation}
All original rows, including zero rows, are counted.
\end{lemma}
\begin{proof}
Independently for every $i$, choose a uniform sign $\sigma_i\in\{-1,1\}$
and a uniform label $a_i\in\{1,2\}$, also independent within each pair.
Set $b_i=(w_i,\sigma_i/\sqrt m)\in\R^{q+1}$ and
$v_i=\sqrt2 e_{a_i}\otimes b_i\in\R^{2(q+1)}$.
The sign removes the off-diagonal covariance with the last coordinate:
\[
 \Exp(b_ib_i^T)=\diag(w_iw_i^T,1/m),\qquad
 \sum_i\Exp(v_iv_i^T)=I_{2(q+1)}.
\]
Each vector has four possible outcomes and squared norm
$2(\norm{w_i}^2+1/m)\leq2\delta$. Theorem~\ref{thm:MSS} therefore supplies
one common outcome with norm at most $(1+\sqrt{2\delta})^2$.
In each label block, the sum is $2\sum_{i:a_i=j}b_ib_i^T$.
Its leading principal block gives the upper Gram bound, and its last scalar
entry gives the upper count bound, both with
$(1+\sqrt{2\delta})^2/2=1/2+e$. Complementing the \emph{same} partition gives
both lower bounds. Signs do not alter the original Grams or cardinalities.
In particular a zero $w_i$ still has its count coordinate. The finite set of
feasible original-channel partitions has a lexicographically first member;
this specifies existence without claiming an efficient search.
\end{proof}

\subsection{Recursive whitening with a noncircular error bound}
Start with $I_0=\{1,\ldots,g\}$ and write
$G_t=\sum_{\alpha\in I_t}u_\alpha u_\alpha^T$, $m_t=|I_t|$.
When $G_t$ is positive, whiten \emph{only for selection}:
$w_\alpha=G_t^{-1/2}u_\alpha$. Apply Lemma~\ref{lem:halving}, choose a
deterministic child $I_{t+1}$, and set
\[
 \delta_t=\max_{\alpha\in I_t}\norm{G_t^{-1/2}u_\alpha}^2+1/m_t,
 \qquad\eta_t=2\sqrt{2\delta_t}+2\delta_t.
\]
Congruence by $G_t^{1/2}$ gives
\begin{equation}\label{eq:recurrence}
 \frac{1-\eta_t}{2}G_t\leq G_{t+1}\leq\frac{1+\eta_t}{2}G_t,
 \quad\frac{1-\eta_t}{2}m_t\leq m_{t+1}\leq\frac{1+\eta_t}{2}m_t.
\end{equation}

\begin{lemma}[Uniform conditioned prefixes]\label{lem:prefix}
Put $C_0=2^{24}$. If $4^k\leq g/(C_0\Lambda q)$, then $2k$ successive
halvings exist, and for every $0\leq t\leq2k$,
\begin{equation}\label{eq:prefix}
 \frac{15}{16}2^{-t}I_q\leq G_t\leq\frac{17}{16}2^{-t}I_q,
 \qquad\frac{15}{16}g2^{-t}\leq m_t\leq\frac{17}{16}g2^{-t}.
\end{equation}
\end{lemma}
\begin{proof}
Assume only the already established bounds at a current stage. Then
\[
 \delta_t\leq\frac{16}{15}2^t(d+1/g)
       \leq\frac{32}{15}2^t d,
\]
because $\Lambda q\geq1$. The current Gram is invertible, justifying that
stage's whitening. Over any justified prefix of $t<2k$, geometric sums give
\begin{equation}\label{eq:error}
 \sum_{t<2k}\eta_t\leq12\sqrt{d4^k}+5d4^k
 \leq\frac{12}{4096}+\frac5{2^{24}}<\frac1{256}.
\end{equation}
Here the displayed full sum denotes a majorant for every justified prefix,
not an assumption that future stages exist. Indeed
$\sum_{t<2k}2^{t/2}=(2^k-1)/(\sqrt2-1)$ and
\[
 \frac{2\sqrt2\sqrt{32/15}}{\sqrt2-1}<12,
 \qquad 2(32/15)<5.
\]
For example $\sqrt2<3/2$, $\sqrt{32/15}<3/2$ and
$(\sqrt2-1)^{-1}=\sqrt2+1<5/2$ already bound the first expression by
$45/4<12$. Also $d4^k\leq2^{-24}$; the last numerator in
\eqref{eq:error} is $49157<65536$ over denominator $2^{24}$.

For nonnegative errors with sum $S<1$,
$\prod(1-\eta_t)\geq1-S$ and
$\prod(1+\eta_t)\leq1/(1-S)$.
Start from $G_0=I_q,m_0=g$. The current loose bounds give the next error
bound; the resulting prefix products lie between $255/256$ and $256/255$,
strictly within $15/16$ and $17/16$. Applying \eqref{eq:recurrence} proves
\eqref{eq:prefix} at the next stage. This completes a prefix induction and
proves existence of all stages, rather than assuming their conditioning.
Every $\eta_t<1$, and both children have positive Gram and positive count.
The case $k=0$ contains no halving and is immediate.
\end{proof}

\subsection{Stopping, original groups and the full row count}
Set
\begin{equation}\label{eq:stop}
 M_* =\min\left\{\frac n{32},\sqrt{\frac g{C_0\Lambda q}}\right\}.
\end{equation}
If $M_*<1$, give \emph{all} channels singleton partitions. This priority
covers $n<32$ or $g<C_0\Lambda q$, including zero frame rows. It will be
counted and checked separately. Otherwise put $k=\lfloor\log_2M_*\rfloor$
and run exactly $2k$ halvings. For $0\leq j<k$ let
$J_j=I_{2j}\setminus I_{2j+2}$, and let $J_k=I_{2k}$ denote the one tail.
These disjoint sets exhaust all original channels. Their \emph{unwhitened}
Grams satisfy
\begin{equation}\label{eq:groupgram}
 G_{J_j}\geq\frac{43}{64}4^{-j}I_q\ (j<k),\qquad
 G_{J_k}\geq\frac{15}{16}4^{-k}I_q.
\end{equation}
The first constant is $(15/16)-(17/64)$, using a lower parent and upper
later subset in their Loewner difference. These are matrix bounds, not traces.

Use the actual spatial pieces of Lemma~\ref{lem:pieces}. Since
$2^k\leq n/32$, their parent lengths satisfy
\[
 L_j\leq\frac{33}{32}n2^{-j}\quad(0\leq j\leq k).
\]
Define $\lambda_j=5L_j^2/(8n^2)$, including the tail. Then
\begin{equation}\label{eq:lambda}
 \lambda_j\leq\frac{5445}{8192}4^{-j}
       <\frac{43}{64}4^{-j},
 \qquad \frac{43}{64}=\frac{5504}{8192}.
\end{equation}
Thus \eqref{eq:groupgram} implies $G_{J_j}\geq\lambda_jI_q$ for every
piece, also the tail. Give channels in $J_j$ the child partition
$\mathcal P_{j+1}$ when $j<k$, and channels in $J_k$ the singleton partition.
Every original channel receives exactly one partition. At $k=0$ all channels
belong to the singleton tail; $G_{J_0}=I_q\geq5I_q/8$ directly.

For $j<k$, $|J_j|\leq(17/16)g4^{-j}$. Hence, with the constant $A_D$ in
Theorem~\ref{thm:fixed},
\[
 N_{\rm coarse}\leq D2^D\frac{17}{16}g
              \sum_{j=0}^{k-1}2^{(D-2)j}
 \leq A_Dg2^{(D-2)k}\leq A_DgV/n^2.
\]
The sum is zero at $k=0$. Since $2^k>M_*/2$,
\[
 |J_k|\leq\frac{17}{16}g4^{-k}
 <\frac{17}{4}\max\{1024g/n^2,C_0\Lambda q\}
 \leq4352g/n^2+\frac{17}{4}C_0\Lambda q.
\]
Multiply by $DV$ for every singleton row to obtain
$N_{\rm out}\leq L_{D,\Lambda}(qV+gV/n^2)$.
In the priority branch the full count $DgV$ is at most
$1024DgV/n^2$ if $n<32$, or $DC_0\Lambda qV$ if $g<C_0\Lambda q$;
the same stated constant pays both. Equalities at either threshold use the
$k\geq0$ branch. No large-size inference is applied at an ineligible input.

\subsection{All-field coercivity with one strong-energy budget}
First take $z$ with componentwise mean zero. Put $f=\mathsf U^Tz$ and
$s=z-\mathsf U f$. These are $q$ weak forms and $g$ strong remainder forms,
respectively; $\mathsf U^Ts=0$. Because $\mathsf U$ is spatially constant,
\begin{equation}\label{eq:split}
 \norm z^2=\norm f^2+\norm s^2,\qquad
 Q_B(z)=a\En(f)+T,\qquad T=H\sum_\alpha\En(s_\alpha).
\end{equation}
For each actual group the parent-cell estimate gives
\[
 H\sum_{\alpha\in J_j}\En(s_\alpha)
 \geq\frac{4H}{L_j^2}\sum_{\alpha\in J_j}\norm{\Delta_js_\alpha}^2.
\]
The channel groups are disjoint: summing these inequalities spends $T$ once.
Since $\Delta_js_\alpha=\Delta_jz_\alpha-u_\alpha\cdot\Delta_j f$, use
\[
 \norm{p-v}^2\geq\frac9{10}\norm v^2-9\norm p^2.
\]
The difference is $10\norm p^2-2\ip p v+\norm v^2/10\geq0$.
Using the actual Grams and \eqref{eq:lambda} yields
\begin{align}\label{eq:strongbudget}
 T&\geq\frac9{10}\sum_{j=0}^k\frac{4H}{L_j^2}
       \ip{\Delta_j f}{G_{J_j}\Delta_j f}
       -36\sum_{j=0}^k\frac H{L_j^2}
                     \sum_{\alpha\in J_j}\norm{\Delta_jz_\alpha}^2\nonumber\\
 &\geq\frac94c\norm f^2-36cb\norm{\mathcal Rz}^2.
\end{align}
The positive pieces sum to $\norm f^2$ by mean zero and their orthogonality.
For the negative term, $\Delta_j$ is a subprojection of its channel's child
observation, or of $I$ at the tail. Also $L_j\geq n2^{-j}$ gives
$H/L_j^2\leq c4^j\leq c4^k$, and stopping implies $4^k\leq b$.
Disjoint channel groups therefore bound the negative sum by the displayed
single observation norm. No factor $q$ or extra copy of strong energy is used.

Independently Lemma~\ref{lem:fourier} gives $T\geq16c\norm s^2$.
Take the convex combination with weights $8/9$ and $1/9$ of this inequality
and \eqref{eq:strongbudget}. It gives
\[
 T+32cb\norm{\mathcal Rz}^2
 \geq2c\norm f^2+\frac{16}{9}c\norm s^2\geq c\norm z^2.
\]
Add $a\En(f)\geq0$. In the priority branch $\mathcal R=I$ and $K=c$
already suffice, so $32cb$ do also.

Finally every channel partition observes every global constant:
$\mathcal R\Pi_0=\Pi_0\mathcal R=\Pi_0$. Its constant and mean-zero images
are orthogonal, $Q_B$ ignores constants, and $32cb\geq c$ pays their full
norm. This proves the all-field statement and, by \eqref{eq:RC}, the precise
normalization in \eqref{eq:fixedupper}. Strong constants on zero weak-frame
rows are included, not assigned a nonexistent spectral gap.

\subsection{Actual-rank lower bound and the cancellation of weak dimension}
Let $C=\bigoplus_\alpha C_\alpha$ be any internally channel-separable linear
map on $\UU$. Let $S_\alpha$ be its row-space orthogonal projection
\emph{inside} $\UU_0$, and $S=\bigoplus_\alpha S_\alpha$. Then
$r=\rank(C|_\UU)=\sum_\alpha\rank S_\alpha$ and $\ker C=\ker S$.
The weak embedding $T_0:\UU_0^q\to\UU$, $T_0f=\mathsf U f$, is an isometry
of real domain dimension $qd_0$. Constant internal coefficients preserve the
actual coexact space. Lemma~\ref{lem:fourier} gives
\begin{equation}\label{eq:weakenergy}
 Q_B(T_0f)=a\En(f)\leq\frac c4\norm f^2.
\end{equation}
Taking a real orthonormal weak basis gives
\begin{equation}\label{eq:ranktrace}
 \tr(T_0^*ST_0)=\sum_\alpha P_{\alpha\alpha}\rank S_\alpha
       \leq\frac{\Lambda q}{g}r.
\end{equation}
Averaging over dimension $qd_0$ therefore supplies a \emph{unit} $f$ with
$t=\norm{ST_0f}^2\leq\Lambda r/(gd_0)$: the two $q$ factors cancel.
Set $w=T_0f$ and $u=(I-S)w$. Then $u\in\UU$, $Cu=0$ and
$\norm u^2=1-t$. Positive-semidefinite polarization gives
\[
 Q_B(u)\leq2Q_B(w)+2Q_B(Sw)\leq c/2+8DHt.
\]
The last step uses $Q_B(v)\leq4DH\norm v^2$ on actual fields.
At zero observation the restoring inequality, even with any finite
$\Phi$ satisfying $\Phi(0)=0$, implies
$c(1-t)\leq c/2+8Dcn^2t$. Thus
$t\geq[2(1+8Dn^2)]^{-1}$, proving the second bound in
\eqref{eq:ranklower} by \eqref{eq:ranktrace}.

Independently $CT_0$ must be injective, since a nonzero weak field in its
kernel would violate restoration by \eqref{eq:weakenergy}. Hence $r\geq qd_0$.
Since $d_0\geq(D-1)V/2$ and
$1+8Dn^2\leq(1+8D)n^2$, their maximum is at least
\[
 \frac{D-1}{8\Lambda(1+8D)}(qV+gV/n^2).
\]
Together with the actual canonical construction this proves the rank order.
Both a minimum on $\UU$ and a minimum requiring all-field restoration obey
these orders, because the lower is already on $\UU$ and the upper is all-field.
There is no output-dimension substitution in the lower argument.

\subsection{Least coefficient for the constructed fixed weights}
For a canonical map, $\mathcal R_\alpha|_{\UU_0}$ is a contraction into the
ambient one-channel field space; its squared singular values sum to at most
its actual rank. That rank equals $\rank(C_\alpha|_{\UU_0})$ by their equal
kernels in \eqref{eq:RC}, even if the projection image leaves $\UU_0$.
Consequently
\[
 \tr(T_0^*\mathcal R^*\mathcal RT_0)
 =\sum_\alpha P_{\alpha\alpha}
             \norm{\mathcal R_\alpha|_{\UU_0}}_{\rm HS}^2
 \leq\frac{\Lambda q}{g}r.
\]
Restoration on weak fields and \eqref{eq:weakenergy} require
$K\norm{\mathcal RT_0f}^2\geq(3c/4)\norm f^2$.
Sum a real orthonormal weak basis and cancel $q$ to obtain
\begin{equation}\label{eq:Ktrace}
 K\geq\frac{3cgd_0}{4\Lambda r}.
\end{equation}
For the constructed map $r\leq L_{D,\Lambda}(qV+gV/n^2)$, so
\[
 K\geq\frac{3(D-1)c}{8\Lambda L_{D,\Lambda}}
          \frac g{q+g/n^2}
 \geq\frac{3(D-1)c}{16\Lambda L_{D,\Lambda}}\min\{g/q,n^2\}.
\]
The last comparison follows by considering which denominator term is larger.
The sufficient all-field $32cb$ proves the opposite bound. A least coefficient
exists: the admissible set is a nonempty closed intersection of restoring
half-lines in $K$, bounded below by \eqref{eq:Ktrace}. This works on $\UU$
or all fields for the same map. These statements complete
Theorem~\ref{thm:fixed}.

\begin{remark}[Endpoints and the role of coherence]\label{rem:endpoints}
For $q=g$, $P=I$ and $M_*<1$; there is no strong complement. The singleton
map has actual rank $gd_0$ and sufficient $K=c$, with $b=1$. The singleton
priority also makes $n=2$, $n<32$ and $g<C_0\Lambda q$ explicit, rather than
an extrapolation of whitening. At $n=32$ or $g=C_0\Lambda q$, $k=0$ is
permitted and the one singleton tail is valid.
Without coherence a coordinate weak $q$-plane can use singleton weak channels
and one root on each strong channel, with actual rank $qd_0$ and $K=c$;
Lemma~\ref{lem:fourier} controls strong mean-zero forms and the roots pay
strong constants. But its diagonal entries equal one, so it is inadmissible
for fixed $\Lambda$ when $g>\Lambda q$. In the remaining range the extra
$gV/n^2$ term is controlled by a $\Lambda$-dependent multiple of $qV$.
This supplies no contrary coherent example.
\end{remark}

\section{Actual checkerboard graphs and plane observations}\label{sec:checkerproof}
We first prove Theorem~\ref{thm:checker}(i) and Theorem~\ref{thm:angle}.
Throughout this section $n$ is even. The parity convention at torus wraps is
part of the family, not an asymptotic simplification.

\subsection{Selected strong cycles and an exact Dirichlet margin}
For the orthogonal checkerboard, channel one has strong outgoing coefficient
$H$ at odd sources and channel two at even sources. Fix one original scalar
spatial/channel component and call its strong-source parity $A$.
Select only its strong $+e_1$ and $+e_D$ edges. Each $A$ vertex has its two
selected outgoing edges, while each opposite-parity vertex has its two
selected incoming edges. The directions are distinct, also at period two.
For an $A$-parity base $u_0$, define
\[
 u_t=u_0+t(e_D-e_1),\qquad v_t=u_t+e_D,\qquad t\in\Z/n\Z.
\]
The translation has order exactly $n$, and parity separates the $u$ and $v$
sets. The selected edges are $u_t$--$v_t$ and $v_t$--$u_{t+1}$, the latter
being the original $+e_1$ edge from $u_{t+1}$. Thus the selected components
are disjoint actual $2n$-cycles and cover all vertices. This is a labelled
finite graph, not a continuum comparison.

Let $\Sigma=\{x:x_D=0\}$. Each cycle has exactly one $u_t\in\Sigma$ and
one $v_{t-1}\in\Sigma$; they are adjacent. Removing these two zero vertices
leaves $2n-2$ free vertices between zero endpoints.
For a path with $N$ free vertices and zero endpoints the sine vectors
$\sin(\pi jt/(N+1))$, $1\leq j\leq N$, diagonalize its Dirichlet matrix
with eigenvalues $4\sin^2(\pi j/[2(N+1)])$.
Indeed the interior second-difference identity gives the eigenvalue and the
endpoint sines vanish; the $N$ eigenvalues are distinct, so the corresponding
nonzero orthogonal vectors span the path space. Its first eigenvalue is at
least $4/(N+1)^2$ by the same sine inequality as before. A path with no free
vertices has zero norm and requires no positive eigenvalue assertion.
Consequently every field zero on $\Sigma$ satisfies
\begin{equation}\label{eq:singleplane-gap}
 Q_B(z)\geq\frac{4H}{(2n-1)^2}\norm z^2
       =(c+\delta_n)\norm z^2,
 \qquad\delta_n=\frac{c(4n-1)}{(2n-1)^2}>0.
\end{equation}
Sum the selected cycle estimates over both original channels and all $D$
spatial components. All other strong edges and all weak terms are merely
discarded as nonnegative. The edge between adjacent zero anchors contributes
zero. At $n=2$ these are actual four-cycles; no parallel-label cancellation
or missing endpoint is used.

In each channel observe every site of $\Sigma$ as a singleton, and the one
remaining rectangle
$\{0,\ldots,n-1\}^{D-1}\times\{1,\ldots,n-1\}$.
It is nonempty because $n\geq2$. This is an exact original-chart partition,
with $n^{D-1}+1$ cells per channel, shared over spatial components. Its block
kernel is zero on $\Sigma$. Lemma~\ref{lem:completion}, with
$M=4DH+c$ and $\delta=\delta_n$, supplies the finite coefficient
$M+2M^2/\delta_n$. Together with \eqref{eq:RC} this proves
Theorem~\ref{thm:checker}(i), including every global constant. This particular
one-plane coefficient is not claimed to have the smaller point/plane-mesh
penalty order in \eqref{eq:pointK}.

\subsection{Every fixed nonzero angle: proof coordinates only}
Let the two weak lines be spanned by unit vectors $v_0,v_1$, with principal
angle $\theta\in(0,\pi/2]$. At parity $\ell$ let $P_\ell=v_\ell v_\ell^T$.
Choose constant unit strong normals $s_\ell\perp v_\ell$ and, for each spatial
component, set $w_\ell=s_\ell\cdot z$. The two-dimensional Gram matrix
$s_0s_0^T+s_1s_1^T$ has eigenvalues $1\pm|s_0\cdot s_1|$, as follows from
its trace two and determinant $1-|s_0\cdot s_1|^2$. Thus
\begin{equation}\label{eq:anglegram}
 \norm{w_0}^2+\norm{w_1}^2\geq\kappa\norm z^2,
 \qquad\kappa=1-|\cos\theta|>0,
\end{equation}
also after summing over sites and spatial components.
These are auxiliary analytic coordinates, not permitted observed rows.
The actual energy contains $H$ times differences of $w_0$ from even sources
and $H$ times differences of $w_1$ from odd sources. All weak terms stay
positive. The selected two-direction cycle proof applies separately to each
scalar functional with its own strong-source parity; these energy contributions
are distinct parts of the actual source-weighted form.

Put $m=\lceil4/\sqrt\kappa\rceil$. If $n<m$, all-channel singletons have
$K=c$ and count $2Dn^D\leq2Dm n^{D-1}$, proving the stated bound directly.
If $n\geq m$, choose the $m$ distinct planes with coordinates
$r_j=\lfloor jn/m\rfloor$, $0\leq j<m$. Their positive cyclic gaps $\ell_j$
are at most $L=\lceil n/m\rceil\leq2n/m$.
Each plane meets each selected cycle in an adjacent pair. Between two
successive anchor pairs there are precisely $2\ell_j-2$ free vertices with
zero endpoints. If $\ell_j=1$ there are none; otherwise the finite Dirichlet
estimate gives eigenvalue at least $4/(2\ell_j-1)^2$.
The wraparound segment is included by the cyclic gap, and zero--zero edges
add zero. For an original field zero on the plane union, both $w_0,w_1$
are zero there, hence
\begin{align}\label{eq:angle-gap}
 Q_B(z)&\geq\frac{4H}{(2L-1)^2}
                      (\norm{w_0}^2+\norm{w_1}^2)\nonumber\\
 &\geq\frac{H\kappa}{L^2}\norm z^2
 \geq\frac{c\kappa m^2}{4}\norm z^2\geq4c\norm z^2.
\end{align}
All inequalities here are non-strict, also at the zero field or when every
vertex is anchored. No nonexistent singleton Dirichlet eigenvalue is used.

In \emph{each original channel}, observe every point of every selected plane
as a singleton. For each nonempty interval between planes use the one full
cross-section rectangle. With $r_0=0$ and $r_m=n$, these intervals are
$\{r_j+1,\ldots,r_{j+1}-1\}$ in the last coordinate; there is no omitted
pre-zero interval. Empty slabs are dropped. This is an exact nonwrapping
original-chart partition with at most $mn^{D-1}+m$ cells per channel.
It is shared across all spatial components and uses only the rows
\eqref{eq:canonical}. A kernel field is pointwise zero on the plane union.
Apply Lemma~\ref{lem:completion} with $\delta=3c$ to
\eqref{eq:angle-gap}; \eqref{eq:pointK} suffices.
This proves Theorem~\ref{thm:angle} without mixing observed channels.

At $\theta=\pi/2$, this is a nonoptimal four-plane variant of the preceding
one-plane construction. At $\theta=0$, $\kappa=0$ and the analytic normal map
is singular; the proof does not apply and does not refute the constant-$P$
theorem. Constants depending on a fixed positive $\theta$ may diverge as the
angle tends to zero. No point-anchor order at a nonorthogonal angle is claimed
by this plane argument.

\section{Distributed point anchors in dimensions three and four}\label{sec:points}
We prove the two smaller full counts in Theorem~\ref{thm:checker}.
The finite point estimate uses all constant modes and counting factors.
Its dimension dependence is classical in character; the exact inequalities
needed here are derived rather than imported from a continuum estimate.

\begin{lemma}[Finite balanced-box point estimates]\label{lem:point}
Let $p=2$ or $3$, let $L\geq1$ be an integer, and let $Q$ be a product of
$p$ integer paths with lengths $\ell_i\in[L/2,L]$. Use counting norm and
internal edge energy $\En_Q$, each path edge once. At any $b\in Q$,
$f(b)=0$ implies
\begin{equation}\label{eq:pointbox}
 \norm f_Q^2\leq
 \begin{cases}
 13L^2(1+\ln L)\En_Q(f),&p=2,\\
 113L^3\En_Q(f),&p=3.
 \end{cases}
\end{equation}
Here $\norm f_Q^2$ denotes $\sum_{y\in Q}|f(y)|^2$.
\end{lemma}
\begin{proof}
Put $N=\prod_i\ell_i$. If $N=1$ the anchored function is zero, so assume
$N>1$. In particular $L\geq2$. The tensor basis from
Lemma~\ref{lem:neumann} has point values and eigenvalues
\[
 |\phi_\nu(b)|^2\leq2^p/N,\qquad
 \lambda_\nu=4\sum_i\sin^2\frac{\pi\nu_i}{2\ell_i}
       \geq\frac{4|\nu|^2}{L^2}\quad(\nu\ne0),
 \qquad0\leq\nu_i<\ell_i.
\]
A singleton factor forces its index to zero. All indices are nonnegative
cosine indices; no signed-frequency multiplicity is suppressed.
On the shell $r=\max_i\nu_i\geq1$, for $p=2$ there are at most
$2r+1\leq3r$ indices, and for $p=3$ at most
$3r^2+3r+1\leq7r^2$. Since $|\nu|^2\geq r^2$ and $r\leq L-1$,
\[
 \sum_{\nu\ne0}|\nu|^{-2}\leq
 \begin{cases}3(1+\ln L),&p=2,\\7L,&p=3.\end{cases}
\]
The first bound follows from
$\sum_{r=1}^{L-1}r^{-1}\leq1+\ln L$ by integration; the second is the
finite sum of at most $L-1$ terms each bounded by seven.
Consequently the nonconstant diagonal inverse satisfies
\begin{equation}\label{eq:finiteGreen}
 G_Q(b,b):=\sum_{\nu\ne0}\frac{|\phi_\nu(b)|^2}{\lambda_\nu}
 \leq\begin{cases}
 3L^2(1+\ln L)/N\leq12(1+\ln L),&p=2,\\
 14L^3/N\leq112,&p=3,
 \end{cases}
\end{equation}
using $N\geq L^p/2^p$.
Expand $f=\sum a_\nu\phi_\nu$ in this counting-orthonormal basis.
Cauchy--Schwarz gives
$|f(b)-\bar f|^2\leq G_Q(b,b)\En_Q(f)$, where $\bar f$ is its arithmetic
mean. The constant-mode mass is $N|\bar f|^2$, not $|\bar f|^2$.
If $f(b)=0$, this mass is at most $12L^2(1+\ln L)\En_Q$ when $p=2$,
and $112L^3\En_Q$ when $p=3$, because $N\leq L^p$.
The positive spectral gap separately bounds the variance by $L^2\En_Q/4$.
Adding the orthogonal constant and variance parts proves \eqref{eq:pointbox}:
$12+1/4<13$, and $112L^3+L^2/4\leq113L^3$.
The $N=1$ case includes $L=1$ and any permissible all-singleton product.
Thus no nonexistent positive gap or lost normalization is used. Point values
were bounded uniformly, so a corner anchor is permitted without an additional
boundary trace assumption.
\end{proof}

\subsection{The full strong graph in original coordinates}
Let $p=D-1$, with $D=3$ or $4$. Fix an original scalar spatial/channel
component and its strong-source parity $A$. For every residue $s$ of that
parity set, with $y\in(\Z/n\Z)^p$,
\begin{equation}\label{eq:paircoordinates}
 x_s(y)=(y_1,\ldots,y_p,s-\textstyle\sum_i y_i)\pmod n,
 \quad u_s(y)=z(x_s(y)),\quad v_s(y)=z(x_s(y)+e_D).
\end{equation}
The first map is a bijection onto sum layer $s$; the shifted map is a
bijection onto layer $s+1$. Even $n$ makes residue parity agree with original
site parity. The $n/2$ pairs $\{s,s+1\}$ cover all sum residues disjointly,
including $\{n-1,0\}$ for odd source parity. Their pair norms therefore sum
to the original scalar counting norm exactly.

The actual strong $+e_D$ edge is $u_s(y)$--$v_s(y)$. For $1\leq i\leq p$,
$x_s(y+e_i)+e_D=x_s(y)+e_i$, so the other strong edges are
$u_s(y)$--$v_s(y+e_i)$. They are \emph{all} $D$ outgoing strong labels,
not a surrogate subset with reassigned endpoint weights. Their exact energy is
\begin{equation}\label{eq:pairenergy}
 E_s=\sum_y|u_s(y)-v_s(y)|^2+
          \sum_{y,i=1}^p|u_s(y)-v_s(y+e_i)|^2.
\end{equation}
A cross edge followed by a rung connects $u_s(y)$ to $u_s(y+e_i)$; all source
translations are generated and every target has a rung, proving connectivity.
At $n=2$ the heads $y,y+e_i$ are distinct. The induced source-torus energy
may repeat an undirected neighbor pair, but every original forward label is
retained. Translations remain bijections.

Insert $v_s(y+e_i)$ between $u_s(y+e_i)$ and $u_s(y)$ and square with factor
two. For the labelled source energy this proves
\begin{equation}\label{eq:graphcomparison}
 \En_p(u_s):=\sum_{y,i=1}^p|u_s(y+e_i)-u_s(y)|^2
 \leq2p\sum_y|u_s-v_s|^2+2\sum_{y,i}|u_s(y)-v_s(y+e_i)|^2
 \leq2pE_s.
\end{equation}
The repeated-rung factor is exactly $2p$, also at period two. Summing $HE_s$
over all layer pairs, both channels and all spatial components gives a lower
part of the original $Q_B$; weak-source terms are only discarded as positive.

\subsection{Balanced quantiles and both layer norms}
For an integer $1\leq m\leq n$, partition each source coordinate into
\[
 I_j=\{\lfloor jn/m\rfloor,\ldots,\lfloor(j+1)n/m\rfloor-1\},
 \qquad0\leq j<m.
\]
These are nonempty paths with lengths $\lfloor n/m\rfloor$ or
$L=\lceil n/m\rceil$. If $L=1$ all are singletons; otherwise a shorter side
is $L-1\geq L/2$. Every product cell therefore meets the exact hypotheses
of Lemma~\ref{lem:point}. Anchor its lowest-coordinate corner as a source
value $u_s(b)=0$, at the actual original point $x_s(b)$.
Cell internal edges are a subset of \eqref{eq:graphcomparison}, counted once;
across-cell and wrap edges are discarded as nonnegative.
Define
\[
 X_2=L^2(1+\ln L),\qquad X_3=L^3.
\]
Lemma~\ref{lem:point} and \eqref{eq:graphcomparison} give
$\norm{u_s}^2\leq52X_2E_s$ for $p=2$ and
$\norm{u_s}^2\leq678X_3E_s$ for $p=3$.
Rungs also give $\norm{v_s}^2\leq2\norm{u_s}^2+2\sum|v_s-u_s|^2$.
As $X_p\geq1$, both original layers satisfy
\begin{equation}\label{eq:pairnorm}
 \norm{u_s}^2+\norm{v_s}^2\leq
 \begin{cases}158X_2E_s,&p=2,\\2036X_3E_s,&p=3.\end{cases}
\end{equation}
Here $3\cdot52+2=158$ and $3\cdot678+2=2036$; no target-layer mass is omitted.
These formulas remain valid at singleton cells because the anchored source
value is zero and the rung still controls its target.

There are $(n/2)m^p$ distinct source anchors per channel: the cell corners
are distinct, \eqref{eq:paircoordinates} is a bijection, and different source
layers are disjoint. Both channels together have $A=nm^p$ channel-labelled
anchors. Use the \emph{same} anchor sites for all $D$ spatial components in
that channel. Although source-coordinate cells need not be original rectangles,
only their original \emph{point} sites are observed. Applying
Lemma~\ref{lem:isolation} independently to the two channels gives actual
canonical partitions with
\begin{equation}\label{eq:pointcount}
 N_{\rm out}\leq2D+2D^2A=2D+2D^2nm^p.
\end{equation}
The initial rows and every channel/component factor are included. Original
chart representatives, not wrapping rectangles or diagonal averages, realize
the points. A block-kernel field is zero at every required anchor; extra
nonpoint means only narrow this kernel.

\subsection{Dimension three: the critical logarithmic mesh}
Take $D=3$, $T_n=1+\ln n$, and let $m$ be the least positive integer with
$m^2\geq4096T_n$. If $m\leq n$, then $L\leq2n/m$ and $L\leq n$, giving
\[
 X_2=L^2(1+\ln L)\leq\frac{4n^2}{m^2}T_n\leq n^2/1024.
\]
Sum \eqref{eq:pairnorm} over the disjoint pairs, channels and components.
Since $H=cn^2$, a field zero at the prescribed anchors satisfies
\begin{equation}\label{eq:D3gap}
 Q_B(z)\geq\frac{H}{158X_2}\norm z^2
          \geq\frac{1024}{158}c\norm z^2\geq4c\norm z^2,
\end{equation}
because $4\cdot158=632<1024$. Put $\rho=\sqrt{4096T_n}\geq1$.
Minimality yields $m=\lceil\rho\rceil\leq2\rho$, so $m^2\leq16384T_n$.
The complete count \eqref{eq:pointcount} is therefore
\[
 N_{\rm out}\leq6+18nm^2\leq6+294912nT_n<2^{19}nT_n.
\]
We used $nT_n\geq1$ and $294912+6<524288$, without divisibility or exact-square
assumptions.

If $m>n$, minimality gives $(m-1)^2<4096T_n$ and the integers satisfy
$n\leq m-1$. The all-channel singleton fallback has
\[
 N_{\rm out}=6n^3<24576nT_n<2^{19}nT_n,\qquad K=c.
\]
This covers every ineligible even $n$, including $n=2$ and $n=m-1$.
At $m=n$ the mesh branch is allowed and its source cells are singleton;
\eqref{eq:D3gap} still follows from the derived rung bound. Thus no guessed
large-size cutoff, empty mesh cell or unproved singleton gap is used.

\subsection{Dimension four: a three-dimensional source mesh}
For $D=4$, first take $n\geq256$ and let $m$ be the least positive integer
with $m^3\geq65536n$. The integer $n$ itself is eligible because
$n^2\geq65536$, so $m\leq n$. Also $L\leq2n/m$ and
\[
 X_3=L^3\leq8n^3/m^3\leq n^2/8192.
\]
The other case of \eqref{eq:pairnorm} now gives, on the actual anchor kernel,
\begin{equation}\label{eq:D4gap}
 Q_B(z)\geq\frac{8192}{2036}c\norm z^2\geq4c\norm z^2,
 \qquad4\cdot2036=8144<8192.
\end{equation}
For $\rho=(65536n)^{1/3}\geq256$, $m=\lceil\rho\rceil\leq2\rho$, hence
$m^3\leq524288n$. Thus $A=nm^3\leq524288n^2$ and
\[
 N_{\rm out}\leq8+32A\leq8+2^{24}n^2\leq2^{25}n^2.
\]
At $n=256$ one has $m=n$ exactly; singleton source cells and their rungs
are valid in \eqref{eq:D4gap}. Noncubes are covered by the same ceiling
argument. For even $n<256$, use all-channel singletons:
$N_{\rm out}=8n^4<524288n^2<2^{25}n^2$ and $K=c$.
No three-dimensional source-cell estimate is applied in that fallback.

\subsection{Finite completion, constants and actual restricted rank}
For either mesh the canonical block kernel obeys
$Q_B(u)\geq4c\norm u^2$ by \eqref{eq:D3gap} or \eqref{eq:D4gap}.
Lemma~\ref{lem:completion} with $\delta=3c$ and $M=4DH+c$ proves the
all-field inequality with exactly \eqref{eq:pointK}. Dividing by $V$ and
using \eqref{eq:RC} gives the stated canonical normalization. On a constant
field $Q_B=0$ and $\mathcal Rz=z$; the finite coefficient is at least $c$,
so these constants are paid. The conservative mesh coefficient can grow as
$cn^4$; in the fallbacks $\mathcal R=I,K=c$ directly.
Restrict only after this all-field proof to the original $\UU_0^2$.
Its actual rank is at most the full count, with no assumed row independence
or coexact invariance of block averaging. This completes
Theorem~\ref{thm:checker}(ii)--(iii).

The same positive stack still has $D^2$ weak and $D^2$ strong eigenvalues per
site. Formula~\eqref{eq:fixedcap} with $q=1,g=2$, including its original
$\alpha_3=729/128$ or $\alpha_4=81/16$, is unchanged and has order $V$.
These improved observations do not reduce the inherited capped budget.
Neither the logarithm in this particular finite construction nor its point
count is asserted necessary for every permitted partition. The present
point proof is for \emph{orthogonal} lines only; it is distinct from the
fixed-angle plane result and supplies no claimed solution to its unproved
point-anchor analogue.

\section{Scope, classical inputs and review context}\label{sec:context}
\subsection{What the comparison establishes}
For fixed $D,\Lambda$, Theorem~\ref{thm:fixed} is uniform in weak dimension
and channel count under a \emph{constant} weak projector. It neither supplies
a basis-free lower nor compares all nonlinear observation classes. Its
finite-support MSS input is nontrivial and classical; the extra signed count
coordinate and the conditioning/count bookkeeping are displayed as an
application, with no historical-priority assertion.
The checkerboard results change only spatial constancy, not local positivity,
local spectra, upper coherence, original channel basis or canonical weights.
They show why that hypothesis cannot be dropped from a per-profile lower of
the same form. Strong-graph geometry, rather than local spectra alone, enters
the constructions. Their improved row upper bounds are not optimal-rank or
uniform-penalty theorems.

The least coefficient assertion for the fixed-space result is deliberately
for the \emph{particular economical map} constructed here. It is not a global
minimum over all canonical maps: all-channel singletons always permit $K=c$
for a positive form, at the cost of many rows. Conversely no practical
conditioning or implementation conclusion follows from finite existence,
lexicographic selection or a rank estimate. The large constants and the
checkerboard coefficients make those limitations substantial.

\subsection{Primary comparisons and their actual extent}
Theorem~\ref{thm:MSS} is the precise published input from
Marcus--Spielman--Srivastava \cite{MSS}. The retained primary reading concerns
Theorem 1.4 and Corollary 1.5 on printed page 3 of arXiv:1306.3969v4,
including the opening corollary argument. It is not a new audit of the full
published MSS proof. The simultaneous count statement used here is proved
with the added coordinate; it is not attributed to an unstated cardinality
clause in that source.

Local coarse averages and their kernels are established mechanisms, as the
selected reading of M\aa lqvist--Peterseim \cite[Sections 2.2--2.3]{MP}
already shows. Finite graph Poincare estimates are classical; the recorded
reading of Levi--Santagati--Tabacco--Vallarino \cite[Section 2, Theorem 2.1]{LSTV}
does not replace the explicit tensor-path proof given here.
Lawler--Limic \cite{LL} provides classical Green-function and potential-kernel
context. The later bounded reading used an author-hosted 289-page draft of
unidentified date, whose pagination differs from the publisher edition:
Section 4.3 and Theorem 4.3.1 (pp.~81--82), Section 4.4.1 and Theorem 4.4.4
(p.~88), and the opening of Section 6.7 (pp.~152--153), with selected
surrounding excerpts, not complete proofs. The first concerns transient
infinite-lattice Green functions in dimensions at least three; the
potential-kernel section supplies two-dimensional logarithmic context.
Boundary rules, generator normalization and removal of the constant mode
must be translated before any constants transfer. Lemma~\ref{lem:point}
derives its finite Neumann constants directly and proves no logarithmic
necessity for all allowed observations.

The GenEO work of Spillane et al.\ \cite{GenEO} is a substantive adaptive
coarse-space predecessor. Operator-dependent local eigenmodes and
partition-of-unity extensions give robust two-level Schwarz convergence.
The retained reading used an author-linked institutional copy of the journal
PDF, including selected setup, Lemma 2.11, Definitions 3.7/3.9,
Assumptions 2.1/3.12/3.13, Lemma 3.21 excerpts, and Theorem 3.22 with
Corollary 3.23 and its size qualification; it was not a full-proof,
implementation or numerical-section reading. These adaptive vector functions
are not prescribed original-channel rectangle means, and their selected
dimension can depend on coefficient geometry and be large. A complete
reduction to our domain, exact row class, uniform rank count and fixed-weight
penalty has not been established; neither has an impossibility of such a
reduction. In particular, failure to establish that implication in a limited
pass provides no evidence of originality. Conversely, our admitted finite
large penalties do not establish the contrast-robust practical conditioning
studied by that solver literature.

Braides--Chiad\`o Piat--Piatnitski \cite{BCP} already studies periodic
high-contrast strong/weak lattice networks. The later bounded reading
includes abstract/introduction excerpts, Section 2 connectivity assumptions
and the statement and preceding decomposition of Theorem 5.1, not its
homogenization proof or flows. A finite paired-layer checkerboard needs an
actual encoding into that theorem's model assumptions before an implication
can be claimed; common terminology is not that encoding. Strong-network
geometry is classical. Implication of the exact finite canonical row bounds
is unresolved here, not proved impossible.
A closer vector/form predecessor needing full comparison is
Budninskiy--Owhadi--Desbrun \cite{BOD}. The institutional primary abstract,
not the full article, was read. It describes operator-adapted differential
forms and differential constraints, so vector forms and operator adaptation
are not themselves available novelty claims. Its complete relation to the
fixed canonical measurements, growing $g,q$, actual restricted ranks and
penalty orders here is unestablished. Failed full-text retrieval is not
evidence of absent overlap.

The new bounded comparison covered four works: MSS, Lawler--Limic, GenEO
and Braides et al.; the other citations have the earlier limited readings
just described. The comparison did not receive or audit this manuscript's
completed source. These are limited primary readings, not an exhaustive literature or
historical-equivalence audit. No external full-text PDF is redistributed in
the source packet. Classical Neumann, Green, PSD trace, strong-graph and
rectangle-isolation methods are acknowledged. The specific rank/penalty
statements require specialist priority scrutiny rather than a claim of a new
general principle.

\subsection*{Model assessments and review context}
At the source-preparation freeze of 11 November 2026, no PDF of this
manuscript had yet been compiled, rendered, read or numerically assessed.
The individual source arguments have retained analytical AI source reports;
those reports are not a whole-manuscript source audit or an exact-PDF review
of this consolidation. A separate integrated-source gate, compilation,
page-by-page rendering review and first exact-PDF assessment are later gates.
There is no numerical grade for this source preparation.

Lluis Eriksson is the named author. Drafting and consolidation assistance was
provided through Codex/OpenAI, exposed as an agent based on GPT-6; the exact
served backend and actual reasoning effort were not exposed and are unknown.
Earlier source audits used separately assigned AI roles in a shared account
and persistent context. No independent, blind, human or fresh-provider review
is established by that role separation. The finite proofs are not Lean or
other formal certificates, and repository checks do not formalize them.

The exact earlier 232B PDF, SHA-256
\begin{center}\small\ttfamily
5595e89c1bcbec1d9b1ae35053a76e6a962f67b45d5d1ac483e5b7a60c9b9be0
\end{center}
received a recorded importance assessment of 4.10/10. That judgment belongs
only to those earlier PDF bytes and supplies no grade here. Its complete
first report, immutable original opinion, same-round protocol clarification
and author response are retained as historical context, without regrading.
The three prior material issues remain unresolved context rather than a new
assessment of this manuscript:
\begin{enumerate}[leftmargin=*]
\item Historical priority and exact implication by predecessors remain
unestablished after limited primary lookups.
\item Practical usefulness, conditioning with allowed penalties, applications
and efficient implementation remain unestablished.
\item The absence of a fatal error in that earlier reading was not
unconditional editorial acceptance. Finite, observation-class, basis/growth
and penalty hypotheses must remain explicit; this consolidation does not
create a new editorial decision or close that issue.
\end{enumerate}
No score transfers from the earlier article or from source reports. The
present scope is finite matrix-gradient observations and the stated actual
flat coexact domain; there is no interacting producer, continuum result or
physical mass-gap conclusion.

The source supplied to a later gated build is
\begin{center}\small\ttfamily\SourceCommit\end{center}
The default \texttt{SOURCE-NOT-BUILT} marks source
preparation only; the build must supply its actual full frozen commit.
Prepared PDF metadata controls are not evidence of bit reproducibility.
That claim requires actual matching PDFs from two fresh clones, which were
not produced in this source-preparation stage.

\appendix
\section{Source correspondence and retained corrections}\label{sec:correspondence}
The source packet retains the complete grouped 235--248 source/audit snapshot,
including rejected variants and adverse reports. Full hashes, exact original
commits and paths are in \texttt{SOURCE-MAPPING.json} and
\texttt{SOURCE-INPUTS.json}; the table gives abbreviated raw SHA-256 identities
for navigation. The manuscript's proofs replace several earlier specialized
selection arguments by the common MSS proof. They do not assert that the new
constructed map is identical to each earlier one.

\begingroup\small
\begin{longtable}{@{}p{.08\textwidth}p{.25\textwidth}p{.36\textwidth}p{.18\textwidth}@{}}
\toprule
Source & Earlier role & Place in this manuscript & Raw hash prefix\\
\midrule\endhead
235 & Dyadic local rectangles, uniform weak line, $g=V$ & Finite Neumann and
piece estimates; special case of Theorem~\ref{thm:fixed} & \texttt{4ccc3829df28}\\
236B & Arbitrary integer sizes; corrected counting normalization &
Lemma~\ref{lem:pieces}, \eqref{eq:RC} and the explicit $VK$ convention & \texttt{b29ea47b3949}\\
237 & Fixed-canonical rank/penalty trace & \eqref{eq:Ktrace}; order recovered
for the new constructed map & \texttt{2d7a29ed2129}\\
238 & Uniform line, $g\geq n^2$ & $q=1$ large-channel case of
Theorem~\ref{thm:fixed} & \texttt{d32783c5ff5d}\\
239 & Uniform line, every $g\geq1$ & $q=1$ rank order; common replacement
construction & \texttt{d277c4773cca}\\
240B & Least coefficient for the specified economical map &
\eqref{eq:penaltyorder}, qualified for the new fixed map & \texttt{d4c37b82e9d9}\\
241B & Balanced rank-one polarization; original cap restored &
$q=1$ scope, \eqref{eq:cap}; no lower row-weight hypothesis here & \texttt{d62cae0af036}\\
242 & Upper-delocalized line without positive minimum row weight &
$q=1$ case; zero rows counted explicitly & \texttt{1016e0c7edc1}\\
243 & Two weak directions, actual matrix groups & $q=2$ case of
Theorem~\ref{thm:fixed}; singleton priority retained & \texttt{985dbb4412fa}\\
244 & Uniform growing weak dimension & Theorem~\ref{thm:fixed},
MSS count augmentation, complete upper and lower proofs & \texttt{59498789a14e}\\
245 & Orthogonal spatial checkerboard &
Theorem~\ref{thm:checker}(i), exact cycle and finite completion & \texttt{6e2aad56076a}\\
246 & Arbitrarily small fixed angle, plane mesh &
Theorem~\ref{thm:angle}; no nonorthogonal point order & \texttt{ab9ab2195aff}\\
247 & Orthogonal dimension four, point anchors &
Theorem~\ref{thm:checker}(iii), \eqref{eq:D4gap} & \texttt{e2ab6a1db681}\\
248 & Orthogonal dimension three, logarithmic point cost &
Theorem~\ref{thm:checker}(ii), \eqref{eq:D3gap} & \texttt{93786fa89262}\\
\bottomrule
\end{longtable}
\endgroup

The original 236A missing-$V$ counting-penalty defect, the 240A small-input
branch defect, and the 241A changed-cap normalization rejection remain bound
to their original sources and reports. Corrected B sources do not turn those
originals into passes. The nonblocking 242 empty-prefix base-case clarification
and 243 singleton-priority clarification are retained too. In this manuscript
the whitening induction has an explicit initial stage and the singleton
priority applies to all channels. These preserved corrections are source
history, not new numerical assessments. The active fixed-angle point-anchor
question has not been developed into a source proof here.

\begin{thebibliography}{9}
\bibitem{MSS}
A.~W.~Marcus, D.~A.~Spielman and N.~Srivastava,
\emph{Interlacing Families II: Mixed Characteristic Polynomials and the
Kadison--Singer Problem}, arXiv:1306.3969v4 (2014), Theorem 1.4.
\url{https://arxiv.org/abs/1306.3969v4}.

\bibitem{LL}
G.~F.~Lawler and V.~Limic, \emph{Random Walk: A Modern Introduction}.
Author-hosted 289-page draft, date unidentified in the inspected version;
pagination differs from the publisher edition.
\url{https://www.math.uchicago.edu/~lawler/srwbook.pdf}.
Selected reading: Section 4.3/Theorem 4.3.1, Section 4.4.1/Theorem 4.4.4,
and the opening of Section 6.7, not complete proofs.

\bibitem{GenEO}
N.~Spillane, V.~Dolean, P.~Hauret, F.~Nataf, C.~Pechstein and R.~Scheichl,
\emph{Abstract robust coarse spaces for systems of PDEs via generalized
eigenproblems in the overlaps}, Numerische Mathematik \textbf{126} (2014),
741--770. \url{https://doi.org/10.1007/s00211-013-0576-y}.
Author-linked institutional journal copy:
\url{https://strathprints.strath.ac.uk/44827/1/10.1007_s00211_013_0576_y.pdf}.
Selected setup, definitions and statements, not the full paper.
\bibitem{MP}
A.~M\aa lqvist and D.~Peterseim,
\emph{Localization of Elliptic Multiscale Problems}, arXiv:1110.0692v4.
\url{https://arxiv.org/abs/1110.0692v4}.
Selected reading: Sections 2.2--2.3.

\bibitem{LSTV}
Levi, Santagati, Tabacco and Vallarino,
\emph{Poincare inequalities on graphs}, arXiv:2203.12530v2.
\url{https://arxiv.org/abs/2203.12530v2}.
Selected reading: Section 2 and Theorem 2.1.

\bibitem{BCP}
Braides, Chiad\`o Piat and Piatnitski,
\emph{Discrete double-porosity models}, arXiv:1406.1752v1 (2014).
\url{https://arxiv.org/abs/1406.1752v1}.
Selected reading: introductory excerpts, Section 2 assumptions and Theorem 5.1 statement/decomposition; not the full proof.

\bibitem{BOD}
M.~Budninskiy, H.~Owhadi and M.~Desbrun,
\emph{Operator-adapted wavelets for finite-element differential forms},
Journal of Computational Physics \textbf{388} (2019), 144--177.
\url{https://doi.org/10.1016/j.jcp.2019.02.018}.
Institutional primary abstract:
\url{https://authors.library.caltech.edu/records/8j84b-1ep03}.
Only that abstract, not the full article, was read in the retained comparison.
\end{thebibliography}

\end{document}
